\chapter{Evaluierung}
\label{chap:Messergebnisse}
Die entwickelten Algorithmen sollen evaluiert werden. Dafür werden zunächst simulative Evaluierung durchgeführt um die Robustheit gegenüber Rauschen und anderem zu überprüfen. Im Anschluss wird das entwickelte System mit einem kommerziell verfügbaren optischen System verglichen. Optische Systeme gelten aktuell als der Goldstandard im Bereich der Bewegungsanalyse.


\section{Verwendete Hardware}
Die Hardware besteht aus drei Komponenten, zum einen ein Auswertecomputer, eine AD/DA-Wandeleinheit und eine Stromverstärkereinheit. Der Auswertecomputer ist ein XXXX. Als AD/DA-Einheit werden zwei Soundkarten des Hersteller $RME\,Audio$ eingesetzt, eine M32 AD Pro und eine M1610 Pro. Diese sind über MADI miteinander verbunden, die Eingangsdaten der M1610 Pro werden dadurch über ein Lichtwellenleiter in Echtzeit an die M32 AD Pro gesendet. Das sorgt dafür, dass die Signale der beiden Einheiten synchronisiert am Auswertecomputer ankommen. Der Auswertecomputer generiert die Anregungssignale für die 3D-Spule und sendet diese an die DA-Einheit. Über die DA-Einheit werden über eine Verstärkereinheit in einen Strom gewandelt. Der Verstärker\footnote{Der Verstärker wurde in Projekt B1/B9 des SFB1261 an der CAU Kiel entwickelt und ein Exemplar für diesen Aufbau überlassen.} ist ein sechskanaliger Verstärker, von dem die ersten drei Kanäle verwendet werden. In Abbildung\,\ref{fig:UsedHardware} ist der Aufbau kurz dargestellt.

\begin{figure}[h!]
  \centering
  \begin{overpic}[width=0.7\textwidth,trim = 0 5cm 0 0]{05_eval/images/UsedHardware.png}
  \end{overpic}
  \caption{Die Abbildung zeigt die verwendete Hardware bestehend aus einem Labornetzteil (rot), einer 32-Kanal-AD-Wandlerkarte (grün), einer 10-Kanal-DA-Wandlerkarte (blau), einer Stromverstärkereinheit (gelb) und eine AVB-Router mit Interface zum Auswertecomputer (pink).}
  \label{fig:UsedHardware}
\end{figure}

Alle Messungen und Simulationen werden mit einem 

\section{Relative Lokalisierung}
\label{sec:rel_lok}




\subsection{Simulative Evaluierung}
In der simulativen Evaluierung sollen 3 unterschiedliche Szenarios untersucht werden, die zunächst kurz definiert werden. Alle Simulationen werden mit einer Abtastrate von $f\tindex{SR}=48000\,\text{kHz}$ und einer Rahmenlänge $i\tindex{Frame} = 512$ durchgeführt. Der Abstand zwischen Quelle und erstem Gelenk wird in allen Simulationen normiert auf 1\,m. Der Asbtand zwischen Sensor und Gelenk wird auf 0.2\,m gesetzt.
\subsubsection{Szenarios und Ergebnisse}
\paragraph{Robustheit gegenüber Rauschen}
Es wird ein Aufbau mit einem Gelenk simuliert, also eine 3D-Spule im Ursprung and dem auch ein kinematisches Element beginnt und ein Gelenk das dieses mit einem weiteren Element verbindet. Die Ansteuerung der Spule erfolgt wie in Kapitel\,\ref{sec:featureExtract} beschrieben. Auf Sensorseite wird WGN (\textit{White Gaussian Noise}) addiert. Diese wird so skaliert, dass das Sensorsignal ein SNR von $10\,\text{dB}$ besitzt. Es wird mit einem Raster von $1^\circ$ der Wertebereich überprüft. Die Simulation wird für 100 Rahmen durchgeführt. Die Bewegung wird an dieser Stelle wie in Kapitel\,\ref{subsubsection:eindeutigkeit} beschrieben, zum einen die Beugung des Gelenks $\phi\tindex{s,2}$ und zum anderen die Richtung in welche das Gelenk gebeugt werden $\phi\tindex{s,1}$.
\begin{figure}[h]
  \centering
  \subfloat[Standardabweichung des MV]{
      \scalefont{2.0}
      \inputTikZ{0.45}{05_eval/images/VarianceMV.tikz}
      \label{fig:VarMV}
  }
  \subfloat[Standardabweichung der Orientierung]{
      \scalefont{2.0}
      \inputTikZ{0.45}{05_eval/images/VarianceProj.tikz}
      \label{fig:VarProj} 
  }
  \caption{Die Abbildung zeigt die Standardabweichung in Abhängigkeit der Gelenksstellung. $\phi\tindex{s,2}$ zeigt dabei wie stark das Gelenk gebeugt ist, $\phi\tindex{s,1}$ zeigt in welche Richtung das Gelenk gebeugt wird.}
  \label{fig:Variance}
\end{figure}

Es wird zunächst geprüft, wie stark der MV und die Orientierung des Elements variieren. Dafür wird zunächst der Durchschnitt von beiden Größen über die 100 simulierten Rahmen gebildet und der Winkel zwischen der Größe und dem Mittelwert gebildet. Diese Winkel werden ausgewertet und die Varianz daraus wird gebildet. Die Standardabweichungen sind in Abbildung\,\ref{fig:Variance} dargestellt. Darin ist die Standardabweichung für jede Haltung dargestellt. Gut zu sehen ist, dass die Varianz des MV relativ konstant ist, mit einem Ausreißer für $\phi\tindex{s,2}\approx 110$. Die Varianz der Orientierung zeigt einen ähnlichen Ausreißer, jedoch zusätzlich eine größere Varianz je kleiner $\phi\tindex{s,2}$ ist.
\begin{figure}[h]
  \centering
  \subfloat[Fehler des Mittelwert]{
      \scalefont{2.0}
      \inputTikZ{0.45}{05_eval/images/ErrorProj.tikz}
      \label{fig:ErrorProj}
  }
  \subfloat[Standardabweichung in Abhängigkeit des Ortes]{
      \scalefont{2.0}
      \inputTikZ{0.45}{05_eval/images/FaktorVar.tikz}
      \label{fig:VarOrt} 
  }
  \caption{ }
  \label{fig:Variance}
\end{figure}
Um diese Ortabhängigkeit zu zeigen wurden die beiden Varianzen zueinander ins Verhältnis gesetzt. Das Ergebnis ist in Abbildung\,\ref{fig:VarOrt} dargestellt. Darin ist gut zu sehen, dass bei gleichbleibender Signalqualität das Ergebnis stärker variiert je kleiner $\phi\tindex{s,2}$ ist. Wichtig ist außerdem, dass der Mittelwert bei genügend großer Signallänge zum richtigen Wert konvergiert, dies wurde in Abbildung\,\ref{fig:ErrorProj} ausgewertet. Der Mittelwert konvergiert in nahezu allen Gelenksstellungen mit dem wahren Wert. Abweichungen gibt es bei $\phi\tindex{s,2}\approx 110$, die bereits eine große Varianz gezeigt haben.

\paragraph{Einfluss von Fehlerfortpflanzung}
Ein Problem bei der Schätzung einer kinematischen Kette ist die Fehlerfortpflanzung. Die Position des ersten Gelenks ist bekannt und wird vor der Laufzeit ausgemessen. Die weiteren Gelenkspositionen werden nach und nach rekonstruiert, schwanken aber je nach Signalqualität. Dieser Effekt soll anhand einer Simulation evaluiert werden. Gegeben sei eine kinematische Kette aus 3 Elementen. 
\begin{figure}[h!]
  \centering
  \begin{overpic}[width=0.7\textwidth,trim = 0 0 0 0]{05_eval/images/FehlerFort.pdf}
  \end{overpic}
  \caption{Der obere Teil der Grafik zeigt, wie die kinematische Kette in der Simulation ausgerichtet ist. Der untere Teil zeigt, wie in der Auswertung das erste kinematische Element angenommen wird. Dieses wird zufällig in einem definiertem Winkelbereich variiert und somit die Position des zweiten Gelenks jeweils leicht falsch bestimmt. Die Orientierung des zweiten Elements wird mit dem beschriebenen Algorithmus bestimmt.}
  \label{fig:FehlerfortDef}
\end{figure}
Die Haltung sei so, dass diese gestreckt sei(Nulllage). Es wird das Sensorsignal für den zweiten Sensor simuliert. Die Orientierungsschätzung des ersten Elements wird zufällig vorgenommen, also die beiden Winkel werden zufällig generiert. Dabei wird die Länge des mittleren Elements $l\tindex{k,1}$(0.1\,m-0.5\,m) variiert, sowie die Amplitude der Winkelvariation $\Delta\phi$ ($0^\circ - 20^\circ$). 
\begin{figure}[h]
  \centering
  \subfloat[Standardabweichung]{
    \scalefont{2}
      \inputTikZ{0.45}{05_eval/images/FehlerfortStdAb.tikz}
      \label{fig:FehlerfortStd}
  }
  \subfloat[Mittlerer Fehler]{
    \scalefont{2}
      \inputTikZ{0.45}{05_eval/images/MeanErrorFort.tikz}
      \label{fig:FehlerfortMaxAb} 
  }
  \caption{Auf der linken Seite ist die Standardabweichung und auf der rechten Seite die maximale Abweichung in $^\circ$, die sich durch die Fehlerfortpflanzung ergeben, in Abhängigkeit von der Länge des mittleren Elements und der Variation dieses.}
  \label{fig:Fehlerfort}
\end{figure}
Das Ergebnis der Simulation ist in Abbildung\,\ref{fig:Fehlerfort} dargestellt. Je größer das mittlere Element ist und je größer die Variation des Winkels, desto größer wird die Standardabweichung. Der Unterschied zwischen Mittelwert über 100 Rahmen und dem wahren Wert ist deutlich kleiner als $1^\circ$.
\paragraph{Einfache Bewegung}
Bewegungen führen zu einer gewissen Unschärfe. Der Algorithmus basiert auf der Detektion von Vorzeichenwechseln, da mindestens zwei Vorzeichenwechsel benötigt werden ist zwischen den $\vec{e}\tindex{zc}$, vergeht zwischen den beiden Nulldurchgängen mindestens eine halbe Periodendauer des Ansteuersignal. Daher gehören die $\vec{e}\tindex{zc}$ zu leicht unterschiedlichen Positionen/Orientierungen. Die Bewegung läuft zyklisch ab. In der Simulation wird $\theta\tindex{Beu} = \omega\,n\,\Delta \,t$ gesetzt, also findet die Bewegung in der $xz$-Ebene statt. $\omega$ wird dabei auf 2\,Hz mit einem Weg von $90^\circ$ gesetzt. Dies entspricht der natürlichen Bewegungsfrequenz\cite{hager-ross_quantifying_2000} und die MCP-Elemente haben einen Bewegungsspielraum von $90^\circ$\cite{lin_modeling_2000}. 

Die Ergebnisse der Simulation sind in Abbildung\,\ref{fig:fig:SimpleMove} dargestellt. Dabei wird ein Fokus auf die $y$-Komponente gelegt. In der Simulation ist die $y$-Komponente zu jedem Zeitpunkt Null, die Schätzung zeigt einen periodischen Verlauf $y$-Komponente mit einer Amplitude von  $0.018$. Die führt zu einer Abweichung der Schätzung von der Bewegungsebene von ca. $1^\circ$.

\begin{figure}[h]
  \centering
  \subfloat[3-Komponenten in Abhängigkeit der Zeit]{
      \scalefont{2}
      \inputTikZ{0.5}{05_eval/images/SimpleMoveVekt.tikz}
      \label{fig:SimpMoveDreiKomp} 
  }
  \subfloat[Y-Komponente]{
      \scalefont{2}
      \inputTikZ{0.5}{05_eval/images/SimpleMoveVektSecond.tikz}
      \label{fig:SimpMoveYKomp} 
  }
  \caption{(a) zeigt den zeitlichen Verlauf der einzelnen Komponenten der Orientierung. (b) ist diesselbe Abbildung mit Fokus auf die $y$-Komponente.} 
   
  \label{fig:SimpleMove}
\end{figure}


\subsubsection{Diskussion}
- warum die Abweichung bei 110?
- Bild wie in Merkmalsextraktion

\subsection{Experimentelle Evaluierung}
Die relative Lokalisierung wurde entwickelt um die Bewegung von Fingern aufzunehmen. Um in einem ersten Test die Algorithmik zu evaluieren, wurde an dieser Stelle ein einfacher Demonstrator bestehend aus drei Elementen, ein längeres Element mit einer 3D-Spule und zwei kürzere mit jeweils einem Sensor. Die Elemente sind durch Gelenke verbunden, die Bewegung in zwei Richtungen ermöglichen, wie bei einem Sattelgelenk. In Abbildung \ref{fig:Demonstrator} ist dieser Aufbau einmal dargestellt.

\begin{figure}[h]
    \centering
    \begin{overpic}[width=0.75\textwidth,trim = 0 0 0 0]{05_eval/images/demonstrator.png}
        \put(5,60){3D-Spule}
        \put(78,32){\parbox{.5in}{Fluxgatemagnetometer}}
        \put(89,6){2\,cm}
        \put(25,25){$\textbf{1}$}
        \put(18,32){$\textbf{2}$}
        \put(6,25){$\textbf{3}$}
        \put(6.5,18){$\textbf{5}$}
        \put(18,15){$\textbf{4}$}
        \put(58,19){$\textbf{1}$}
        \put(48,19){$\textbf{2}$}
        \put(60.5,4){$\textbf{3}$}
        \put(50,4){$\textbf{4}$}
        \put(66,49){$\textcolor{blue}{z}$}
        \put(86,44){$\textcolor{green}{y}$}
    \end{overpic}
    \caption{Demonstrator: Die Abbildung zeigt den Demonstrator, der für die Evaluierung verwendet wird. Die 3D-Spule und die kinematischen Elemente sind mit optischen Markern ausgerüstet.  }
    \label{fig:Demonstrator}
\end{figure}

Zum Vergleich wird ein optisches System der Firma $Qualisys AB$ verwendet, bestehend aus 8 Miqus Kameras. Diese sind um den Demonstrator herum aufgebaut. Für die optische Bewegungsaufnahme wurden optische Marker zum einen auf der 3D-Spule und zum anderen auf den kinematischen Elementen befestigt, diese werden verwendet um die Orientierung der Elemente aufzunehmen und werden als $Ground Truth$ verwendet. Der Aufbau ist in Abbildung\,\ref{fig:aufbauEval} zu sehen.

\begin{figure}[h]
  \centering
  \begin{overpic}[width=0.75\textwidth,trim = 0 0 0 0]{05_eval/images/EvalAufbau_Pic.png}
  \end{overpic}
  \caption{Evaluierungsaufbau: Das Foto zeigt den Demonstrator in der Mitte (rot) und im Außenbereich die Kameras (gelb).}
  \label{fig:aufbauEval}
\end{figure}

Es wurden zum einen statische und dynamische Messungen durchgeführt. Die Länge der Aufnahmen waren immer 15\,s. 

\subsubsection{Räumliche Synchronisierung und Schätzung der räumlichen Größen}

Die beiden Systeme arbeiten in unterschiedlichen Koordinatensystemen. Das magnetische System spannt ein Koordinatensystem basierend auf der Ausrichtung der 3D-Spule auf, also $x$-Richtung = Ausrichtung der X-Spule\dots . Das Koordinatensystem des optischen Systems wird in einer Kalibrierungsprozedur definiert, durch eine festgelegte Markeranordnung. Die Definition der beiden Koordinatensysteme ist in Abbildung \ref{fig:CoordSystems}.

\begin{figure}
    \centering
    \subfloat[Magnetisches Koordinatensystem]{
        \begin{overpic}
            [width=0.5\textwidth,trim = 0 0 0 0]{05_eval/images/Coordinate_transformation.png}
        \end{overpic}
        \label{fig:MagCoord}
    }
    \subfloat[Optisches Koordinatensystem]
    {
        \begin{overpic}
            [width=0.5\textwidth,trim = 0 0 0 0]{05_eval/images/OptischesKoord.png}
        \end{overpic}
        \label{fig:OptCoord}
    }
    \caption{Definition der Koordinatensysteme: Der rote Pfeil ist jeweils die $x$-Achse, grün die $y$-Achse und blau die $z$-Achse. In (a) ist eine 3D-Spule mit einem 3D-Druckaufsatz, an dem optische Marker befestigt sind. Mit den farbigen Pfeilen ist das lokale Koordinatensystem der Spule dargestellt. In (b) ist das Kalibrationswerkzeug für das optische System. Während der Kalibrierung wird die Platte mit den Markern in der Mitte des Messvolumens platziert. Währenddessen wird der Stab mit den zwei Markern durch den Messbereich bewegt. Die Lage und Orientierung der Platte während der Kalibrierung definiert dann die Orientierung des optischen Koordinatensystem.}
    \label{fig:CoordSystems}
\end{figure}

Um die beiden Ergebnisse miteinander vergleichen zu können, müssen diese im gleichen Koordinatensystem sein. Ein Problem ist, dass der Demonstrator aus Abbildung \ref{fig:Demonstrator} nur schwer so positionert werden kann, dass beide Koordinatensysteme miteinander ausgerichtet sind. Daher sind an der 3D-Spule optische Marker angebracht, mit deren Hilfe die magnetischen Koordinatenachsen im optischen Koordinatensytem detektiert werden können. Aus den so detektierten Einheitsvektoren ($\vec{e}\tindex{x,Mag},\vec{e}\tindex{y,Mag},\vec{e}\tindex{z,Mag}$) ergibt sich eine Rotationsmatrix $\bvec{R}\tindex{Opt$\rightarrow$ Mag}$ zwischen den beiden Koordinatensystemen 

\begin{equation}
    \bvec{R}\tindex{Opt$\rightarrow$ Mag} = \left[\begin{array}{c} \vec{e}\tindex{x,Mag}^\text{\,T} \\ \vec{e}\tindex{y,Mag}^\text{\,T} \\ \vec{e}\tindex{z,Mag}^\text{\,T} \end{array}\right].
\end{equation}

Die optische Orientierung der kinematischen Elemente wird mithilfe der Positionen $\vec{r}\tindex{$i,j$}$ der optischen Marker bestimmt, wobei $i$ die Nummer des Elements ist und $j$ die Nummer des jeweiligen Markers ist. Die Orientierungsschätzung $\hat{\vec{e}}_{\text{opt},i}$ des jeweiligen Elements ergibt sich wie folgt
\begin{equation}
    \hat{\vec{e}}_{\text{opt},i} = \frac{(\vec{r}_{i,1} - \vec{r}_{i,2}) + (\vec{r}_{i,3} - \vec{r}_{i,4}) }{ |(\vec{r}_{i,1} - \vec{r}_{i,2}) + (\vec{r}_{i,3} - \vec{r}_{i,4})|  }.
\end{equation}
Durch Anwendung von $\bvec{R}\tindex{Opt$\rightarrow$ Mag}$ wird diese in das magnetische Koordinatensystem transformiert

\begin{equation}
    \hat{\vec{e}}\tindex{opt,$i$,mag} =  \bvec{R}\tindex{Opt$\rightarrow$ Mag}\,\hat{\vec{e}}_{\text{opt},i},
\end{equation}
und kann mit der magnetischen Schätzung verglichen werden.
\subsubsection{Kalibrierung}
Der Algorithmus für die Schätzung einer kinematischen Kette benötigt als Vorwissen die Abmessungen der kinematischen Kette. Die benötigten Parameter sind die zum einen die Position des ersten Gelenks $\vec{r}\tindex{j}$, sowie für jedes Element der Abstand von Gelenk zu Sensor $l\tindex{js,$i$}$, die Länge der Elemente $l\tindex{ke,$i$}$ und die Dicke der Elemente $d\tindex{ke,$i$}$.

\begin{figure}
    \centering
    \begin{overpic}[width=0.75\textwidth]{05_eval/images/Kalibrierung.pdf}
      \put(20,17){$r\tindex{j}$}
      \put(45,25){$l\tindex{js,1}$}
      \put(60,17){$d\tindex{ke,1}$}
      \put(68,28){$l\tindex{ke,1}$}
      \put(60,45){$l\tindex{js,2}$}
      \put(75,42){$d\tindex{ke,2}$}
      \put(58,10){$x$}
      \put(27,48){$y$}
      \put(12,60){$z$}
    \end{overpic}
    \caption{Kalibrierung der kinematischen Kette: Die Abbildung zeigt eine kinematische Kette bestehend aus drei Elementen. Die Roten Pfeile zeigen alle Größen die für den Algorithmus vor der Laufzeit bestimmt werden müssen.}
    \label{fig:TBCalibrated}
\end{figure}
Für die Kalibrierung wurden zunächst alle Größen händisch mit einem Lineal ausgemessen. Zwei Schwierigkeiten dabei sind, dass es nicht möglich ist von der Mitte der Quelle aus zu messen und die Länge des Sensorelements beträgt ca. 2\,cm, hier ist es unklar an welcher Stelle genau die Position für einen modellhaften Punktsensor optimal gelegt werden soll. Die gemessenen Initialwerte waren $\vec{r}\tindex{j} = (-0.02,0.00,0.29)\,\text{m}$, $d\tindex{ke,1}=0.015\,\text{m}$, $l\tindex{js,1}=0.065\,\text{m}$, $d\tindex{ke,2}=0.015\,\text{m}$, $l\tindex{js,2}=0.065\,\text{m}$.

Daher wurden die statischen Messungen verwendet um die Kalibrierung durchzuführen. Begonnen wurde mit dem ersten Element, dafür wurden 13 Aufnahmen in unterschiedlichen Ausrichtungen verwendet. Die optimale Kalibrierung sei definiert, dass für diese die optische $\hat{\vec{e}}\tindex{opt,mag}^{\,(j)}$ und die magnetische $\hat{\vec{e}}\tindex{mag}^{\,(j)}$ Schätzung am besten miteinander übereinstimmen. Als Differenz zwischen den jeweiligen Schätzungen wird der Winkel zwischen diesen verwendet. Also Kostenfunktion ergibt sich daraus dann die folgende Gleichung, wobei $N$ die Anzahl der Messungen und $j$ die Nummer des Elements ist

\begin{equation}
    e = \frac{1}{N} \sum_{j=1}^{N} \arccos(\hat{\vec{e}}\tindex{opt,mag}^{\,(j)}\cdot\hat{\vec{e}}\tindex{mag}^{\,(j)}).
        \label{eq:CostFkt}
\end{equation}
In einem Bereich von $\pm 2\,cm$ zu den gemessenen Werten wurden in einem Raster von 1\,mm untersucht und die Kombination mit dem geringsten Fehler wurde übernommen. Mit den Werten des ersten Elements wurde dann das zweite Element auf die selbe Art kalibriert, hier mit $N=8$ Aufnahmen. Die Kalibriertenwerte unterscheiden sich leicht zu den händisch gemessenen Werten und lauten wie folgt $\vec{r}_j = (-0.0137,0.0029,0.2960)\,\text{m}$, $d\tindex{ke,1}=0.008\,\text{m}$, $l\tindex{js,1}=0.063\,\text{m}$, $d\tindex{ke,2}=0.008\,\text{m}$, $l\tindex{js,2}=0.063\,\text{m}$.

\subsubsection{Ergebnisse}
In den Tabellen\,\ref{tab:Ergebnisse1},\,\ref{tab:Ergebnisse2} sind die Ergebnisse der statischen Messungen zu sehen. Darin sind die Gelenkswinkel der magnetischen und der optischen Schätzung aufgetragen. Der mittlere Fehler (\textit{Mean Absolute Error}) für das erste Element beträgt $1.3^\circ$ und für das zweite Element $2.7^\circ$.
\begin{table}[h]
    \centering
    \begin{tabular}{|c| c| c| c| c|}
    \hline 
     $\hat{\theta}_{opt,1}$ & $\hat{\theta}_{mag,1}$ & $\hat{\phi}_{opt,1}$ & $\hat{\phi}_{mag,1}$ & Fehler \\ 
     \hline
     $-3.5^\circ$ & $-3.7^\circ$ & $1.2^\circ$  & $1.7^\circ$ & $0.5^\circ$ \\ 
     \hline
     $-3.6^\circ$ & $-3.8^\circ$ & $42.3^\circ$   & $42.1^\circ$ & $0.3^\circ$ \\
     \hline
     $-3.9^\circ$ & $-4.4^\circ$ & $88.5^\circ$  & $86.6^\circ$ & $1.9^\circ$ \\
     \hline
     $-1.3^\circ$ & $-1.9^\circ$ & $-39.2^\circ$ & $-39.6^\circ$ & $0.7^\circ$ \\
     \hline
     $-0.2^\circ$ & $-0.1^\circ$ & $-88.0^\circ$ & $-88.4^\circ$ & $0.4^\circ$ \\ 
     \hline
     $-45.6^\circ$ & $-46.0^\circ$ & $5.9^\circ$ & $3.1^\circ$ & $2.0^\circ$  \\ 
     \hline
     $-45.4^\circ$ & $-45.3^\circ$ & $44.1^\circ$  & $42.9^\circ$ & $0.8^\circ$ \\
     \hline
     $-46.5^\circ$ & $-46.3^\circ$ & $-32.3^\circ$ & $-35.6^\circ$ & $2.3^\circ$ \\
     \hline
     $35.0^\circ$ & $34.9^\circ$ & $2.6^\circ$ & $3.9^\circ$ & $1.1^\circ$ \\
     \hline
     $35.6^\circ$ & $34.7^\circ$ & $41.2^\circ$  & $41.2^\circ$ & $0.8^\circ$ \\ 
     \hline
     $34.7^\circ$ & $35.4^\circ$ & $-48.7^\circ$  & $-47.7^\circ$ & $1.1^\circ $ \\ 
     \hline
     $86.6^\circ$ & $85.8^\circ$ & $171.2^\circ$ & $132.6^\circ$ & $2.6^\circ$ \\
     \hline
     $-89.4^\circ$ & $-87.1^\circ$ & $-14.5^\circ$ & $-57.6^\circ$ & $2.4^\circ$ \\
     \hline
    \end{tabular}
    \caption{Ergebnisse der statischen Messung des ersten kinematischen Elements.}
    \label{tab:Ergebnisse1}
  \end{table}

  \begin{table}[h]
    \centering
    \begin{tabular}{|c| c| c| c| c|} 
      \hline
      $\hat{\theta}_{\text{opt},2}$ & $\hat{\theta}_{\text{mag},2}$ & $\hat{\phi}_{\text{opt},2}$ & $\hat{\phi}_{\text{mag},2}$ & Fehler \\ 
     \hline
     $3.7^\circ$ & $3.8^\circ$ & $89.3^\circ$ & $86.3^\circ$ & $3.0^\circ$ \\ 
     \hline
     $0.0^\circ$ & $2.3^\circ$ & $-89.2^\circ$ & $-89.3^\circ$ & $2.2^\circ$ \\
     \hline
     $-1.1^\circ$ & $1.1^\circ$ & $-54.6^\circ$ & $-54.9^\circ$ & $1.9^\circ$ \\
     \hline
     $1.6^\circ$ & $3.7^\circ$  & $40.2^\circ$ & $40.0^\circ$ & $2.0^\circ$ \\
     \hline
     $86.8^\circ$ & $84.1^\circ$ & $6.1^\circ$ & $40.3^\circ$ & $2.8^\circ$ \\ 
     \hline
     $32.1^\circ$ & $34.8^\circ$ & $1.7^\circ$ & $3.3^\circ$ & $2.7^\circ$ \\ 
     \hline
     $-35.8^\circ$ & $-33.8^\circ$ & $2.5^\circ$ & $0.5^\circ$ & $2.1^\circ$ \\
     \hline
     $-85.5^\circ$ & $-86.5^\circ$ & $165.4^\circ$ & $-116.0^\circ$ & $4.1^\circ$ \\
     \hline
    \end{tabular}
    \caption{Ergebnisse der statischen Messung des zweiten kinematischen Elements.}
    \label{tab:Ergebnisse2}
  \end{table}
In Abbildung\,\ref{fig:RawResults} sind die Ergebnisse der dynamischen Aufnahme ohne Nachverarbeitung dargestellt. Das Kamerasystem liefert mit einer Abtastrate von 100\,Hz Positionen. Um beide System gut miteinander vergleichen zu können, wurde die Rahmenlänge der Signalverarbeitung so angepasst, dass auch die magnetische Schätzung Werte mit 100\,Hz liefert.
  \begin{figure}[h]
    \centering
    \subfloat[Gelenkwinkel für das erste Gelenk]{
      \inputTikZ{0.425}{05_eval/images/FirstElementAnglesRaw.tikz}
      \label{fig:AngleFirstraw}
    }
    \subfloat[Gelenkwinkel für das zweite Gelenk]{
      \inputTikZ{0.425}{05_eval/images/SecondElementAngleRaw.tikz}
      \label{fig:AngleSecondraw}
    }
    \caption{Die Abbildung zeigt die beiden Gelenkswinkel für das erste(a) und zweite Gelenk(b). Dabei sind die Schätzungen der magnetischen Lokalisierung ohne Nachverarbeitung(gestrichelt), sowie die optischen Schätzungen über der Zeit aufgetragen. }
    \label{fig:RawResults}
  \end{figure}

  \begin{figure}[h]
    \centering
    \subfloat[Gefilterte Winkel der Bewegung für das erste Element]{
      \inputTikZ{0.425}{05_eval/images/FirstElement_Angle}
      \label{fig:AngleFirst}
    }
    \subfloat[Absoluter Fehler zwischen optischer und magnetischer Schätzung]{
      \inputTikZ{0.425}{05_eval/images/FirstElement_Error}
      \label{fig:ErrorFirst}
    }
    \caption{Ergebnisse des ersten Elements: (a) zeigt die Winkelschätzungen im Vergleich. Die optischen sind in gestrichelten Linien und die magnetischen mit durchgezogenen Linien dargestellt. (b) zeigt den Winkelfehler zwischen den beiden SChätzungen.}
    \label{fig:ResultsFirstElement}
  \end{figure}
  
  
  \begin{figure}[h]
    \centering
    \subfloat[Gefilterte Winkel der Bewegung für das zweite Element]{
      \inputTikZ{0.425}{05_eval/images/SecondElement_Angle}
      \label{fig:AngleSecond} 
    }
    \subfloat[Absoluter Fehler zwischen optischer und magnetischer Schätzung]{
      \inputTikZ{0.425}{05_eval/images/SecondElement_Error}
      \label{fig:ErrorSecond}
    }
    \caption{Ergebnisse für das zweite Element: Die Darstellung ist wie in Abbildung\,\ref{fig:ResultsFirstElement}.}
    \label{fig:ResultsSecondElement}
  \end{figure}
\subsubsection{Diskussion}

\section{Absolute Lokalisierung}
\begin{figure}[h]
  \centering
  \subfloat[Simulationssetup]{
  \scalefont{2}
  \inputTikZ{0.6}{05_eval/images/AbsSimSensorSetup.tikz}
  \label{fig:SimAbsSetup}
  }
  \subfloat[Laboraufbau]{
    \begin{overpic}[width=0.4\textwidth,trim = 0 0 0 0]{05_eval/images/AufbauAbslocal.png}
    \end{overpic}
    \label{fig:SimAbsSetup}
  }
  \caption{Setup}
  \label{fig:AbsSetup}
\end{figure} 
\subsection{Fehlermaß}
Bei der Auswertung der absoluten Lokalisierung müssen zum einen die Position und die Orientierung evaluiert werden. Um eine Evaluierung durchführen zu können, muss ein Fehlermaß eingeführt werden. Bei der Position ist dieses naheliegend die euklidische Abstand zwischen der magnetischen Schätzung und dem Ergebnis der optischen Schätzung. Beim Vergleich der Orientierung ist dies komplexer, da bei Eulerwinkeln keine euklidische Distanz existiert. Hier soll ein Skalar eingeführt werden, der ein Maß für die Qualität der Schätzung ist. 

Die Orientierung der 3D-Spule $\bvec{R}\tindex{RPY}$ ist die Verdrehung des Spulekoordinatensystem im Vergleich zum globalen Koordinatensystem. Die einzelnen Achsen des lokalen Koordinatensystems beschrieben im globalen Koordinatensystem sind die Spalten der Rotationsmatrix

\begin{equation}
    \bvec{R}\tindex{RPY} = \begin{pmatrix}
        \vec{e}_x^{\,\text{glob}} & \vec{e}_y^{\,\text{glob}} & \vec{e}_z^{\,\text{glob}}.
    \end{pmatrix}
\end{equation}

In der Evaluierung werden diese magnetisch und optische geschätzt, also die Koordinatenachsen der magnetischen und optischen Schätzung. Als Maß wird der absolute Winkel verwendet um den die beiden Koordinatensysteme zueinander verdreht sind. Dafür wird zunächst Die Rotationsmatrix $\bvec{R}\tindex{Eval}$ zwischen den beiden Koordinatensystemen bestimmt

\begin{equation}
    \bvec{R}\tindex{Eval} = \bvec{R}\tindex{RPY,mag} (\bvec{R}\tindex{RPY,opt})^{-1}.
\end{equation}
Aus $\bvec{R}\tindex{Eval}$ wird ein äquivalentes Quaternion bestimmt. Quaternionen sind Vektoren mit vier Komponenten, die in der Informatik häufig eingesetzt werden um räumliche Rotationen zu beschreiben. Der Vorteil ist, dass diese weniger Speicherplatz benötigen Referenz. Die Komponenten eines Quaternions $\bvec{q}$ kann mit der Rotationsachse $\vec{e}\tindex{Rot}$ und dem Drehwinkel $\phi$ beschrieben werden

\begin{equation}
    \bvec{q} = \begin{pmatrix}
        \cos(\frac{\phi}{2}) \\
        \sin(\frac{\phi}{2})\,e\tindex{Rot,x} \\
        \sin(\frac{\phi}{2})\,e\tindex{Rot,y} \\
        \sin(\frac{\phi}{2})\,e\tindex{Rot,z} 
    \end{pmatrix} = \begin{pmatrix}
      w \\
      x \\
      y \\
      z 
  \end{pmatrix} .
\end{equation}

Aus $\bvec{R}\tindex{Eval}$ wird das zugehörige Quaternion $\bvec{q}\tindex{Eval}$ bestimmt, wobei hier nur die erste Komponente $w$ benötigt wird. Zunächst werden die drei Eulerwinkel $\theta\tindex{Eval},\phi\tindex{Eval},\psi\tindex{Eval}$ wie in den Gleichungen\,\ref{eq:ThetaOri}\,,\ref{eq:PsiOri}\,,\ref{eq:PhiOri} aus $\bvec{R}\tindex{Eval}$ bestimmt. Daraus ergeben sich drei Quaternionen $\bvec{q}\tindex{z},\bvec{q}\tindex{y},\bvec{q}\tindex{x}$

\begin{equation}
  \bvec{q}\tindex{z} = \begin{pmatrix}
    \cos(\frac{\psi\tindex{Eval}}{2}) \\
    \sin(\frac{\psi\tindex{Eval}}{2})\,0 \\
    \sin(\frac{\psi\tindex{Eval}}{2})\,0 \\
    \sin(\frac{\psi\tindex{Eval}}{2})\,1 
\end{pmatrix},
\label{eq:quatZ}
\end{equation}

\begin{equation}
  \bvec{q}\tindex{y} = \begin{pmatrix}
    \cos(\frac{\theta\tindex{Eval}}{2}) \\
    \sin(\frac{\theta\tindex{Eval}}{2})\,0 \\
    \sin(\frac{\theta\tindex{Eval}}{2})\,1 \\
    \sin(\frac{\theta\tindex{Eval}}{2})\,0 
\end{pmatrix},
\label{eq:quatY}
\end{equation}
\begin{equation}
  \bvec{q}\tindex{x} = \begin{pmatrix}
    \cos(\frac{\phi\tindex{Eval}}{2}) \\
    \sin(\frac{\phi\tindex{Eval}}{2})\,1 \\
    \sin(\frac{\phi\tindex{Eval}}{2})\,0 \\
    \sin(\frac{\phi\tindex{Eval}}{2})\,0 
\end{pmatrix}.
\label{eq:quatX}
\end{equation}
Die drei einzelnen Quaternionen werden aufmultipliziert und man erhält
\begin{equation}
  \bvec{q}\tindex{Eval} = \bvec{q}\tindex{z} \circ \bvec{q}\tindex{y} \circ \bvec{q}\tindex{x}.
  \label{eq:QuatEval}
\end{equation}
Die Komponenten zwei bis vier von $\bvec{q}\tindex{Eval}$ werden verworfen und aus der ersten Komponente wird die effektive Verdrehung $\alpha\tindex{Eval}$ bestimmt
\begin{equation}
  \alpha\tindex{Eval} = 2\,\arccos(q\tindex{Eval,w}).
  \label{eq:EffVerd}
\end{equation}
Dieser Winkel wird als Fehlermaß für die Orientierungsschätzung verwendet. 
 

\subsection{Simulative Evaluierung}


\paragraph{Robustheit der Positionsschätzung gegen Rauschen}
Die Robustheit der Positionsschätzung soll anhand eines Sensorpaares, also in einer Koordinate, geprüft werden. Da die drei Koordinaten unabhängig voneinander bestimmt werden, ergibt sich jeweils dasselbe Ergebnis. Auf das Sensorsignal wird weißes Rauschen addiert, mit einem konstatnten Rauschteppich von $150\,\frac{\text{pT}}{\sqrt{\text{Hz}}}$ ergibt. Die Spule wird so angesteuert, dass der maximale Strom bei 1\,A liegt. 
\begin{figure}[h]
  \centering
  \subfloat[Standardabweichung der Position]{
    \scalefont{2}
    \inputTikZ{0.425}{05_eval/images/StdAbweichung_Pos.tikz}
    \label{fig:stdAbwPosSim}
  }
  \subfloat[Abweichung der mittleren Position von Ground Truth]{
    \scalefont{2}
    \inputTikZ{0.425}{05_eval/images/absoluterFehler_Pos.tikz}
    \label{fig:errorPosSim}
  } 
  \caption{}
  \label{fig:PosSimEval}
\end{figure}

\paragraph{Einfluss eines Positionsfehler auf die Orientierungsschätzung}
\begin{figure}[h]
  \centering
  \scalefont{2}
  \inputTikZ{0.425}{05_eval/images/Abhängigkeit_or_Pos.tikz}
  \label{fig:stdAbwOrPosVar}

  \caption{}
  \label{fig:PosSimEval}
\end{figure}
\paragraph{Einfache Bewegung}
\begin{figure}[h]
  \centering
  \subfloat[Angle of movement for the second element]{
    \scalefont{2}
    \inputTikZ{0.45}{05_eval/images/MovementComponent.tikz}
    \label{fig:MovementComponent}
  }
  \subfloat[Magnitude of the angular error between optical and magnetic estimation]{
    \scalefont{2}
    \inputTikZ{0.45}{05_eval/images/FehlerBewegung.tikz}
    \label{fig:ErrorMovement}
  } 
  \caption{}
  \label{fig:AbsMovement}
\end{figure}

\subsection{Experimentelle Evaluierung}
\subsubsection{Räumliche Synchronisierung}

\begin{figure}[h]
  \centering
  \begin{overpic}[width=0.7\textwidth,trim = 0 0 0 0]{05_eval/images/SpatialSync_Abslocal.png}
  \end{overpic}
  \caption{
  }
  \label{fig:SpatialSyncAbslocal}
\end{figure}

\subsubsection{Ergebnisse}
\begin{figure}[h]
  \centering
  \scalefont{2}
  \inputTikZ{0.45}{05_eval/images/DynEvalCompsPos.tikz}
  \label{fig:stdAbwOrPosVar}
\end{figure}

\begin{figure}[h]
  \centering
  \subfloat[$\theta$]{
    \scalefont{2}
    \inputTikZ{0.45}{05_eval/images/DynEvalAngleTheta.tikz}
    \label{fig:MovementComponent}
  }
  \\
  \subfloat[$\phi$]{
    \scalefont{2}
    \inputTikZ{0.45}{05_eval/images/DynEvalAnglePhi.tikz}
    \label{fig:ErrorMovement}
  } 
  \\
  \subfloat[$\psi$]{
    \scalefont{2}
    \inputTikZ{0.45}{05_eval/images/DynEvalAnglePsi.tikz}
    \label{fig:ErrorMovement}
  } 
  \caption{}
  \label{fig:AbsMovement}
\end{figure}

\begin{figure}[h]
  \centering
  \subfloat[$\theta$]{
    \scalefont{2}
    \inputTikZ{0.45}{05_eval/images/DynEvalPosFehler.tikz}
    \label{fig:MovementComponent}
  }
  \\
  \subfloat[$\phi$]{
    \scalefont{2}
    \inputTikZ{0.45}{05_eval/images/DynEvalWinkelfehler.tikz}
    \label{fig:ErrorMovement}
  } 
  \caption{}
  \label{fig:AbsMovement}
\end{figure}
  
\subsection{Diskussion}


